% DRAFT — not for submission. Numbers come from results/phase3_*; regenerate with scripts/analyze_calibration.py.
% Every reference marked [to verify] must be read in full before submission.
\documentclass[11pt]{article}
\usepackage[margin=1in]{geometry}
\usepackage{booktabs,hyperref}
\title{Ten Minutes of Shared Listening: Label-Free Functional Alignment\\ for Timing-Free MEG Word Decoding in New Subjects}
\author{Joaquim Manji-Jung\\ Paris}
\date{Draft, November 2026}
\begin{document}
\maketitle

\begin{abstract}
Non-invasive decoders of heard words are usually trained and tested on the same person, and recent work showed
that much of the reported accuracy of contextual decoders comes from word timing leaked by overlapping input
windows rather than from brain activity \cite{jayalath2026shortcuts}. We study the setting that matters for
patients: a \emph{new} person, a decoder free of the timing shortcut, and a few minutes of calibration.
On LibriBrain100 \cite{libribrain100}, a decoder trained on one subject and applied one word window at a time
transfers weakly to 28 never-seen subjects (24.9\% balanced top-10 accuracy over 50 words, chance 20\%).
Supervised fine-tuning on the new subject's first minutes helps (+2.7 points with 5 minutes, +6.3 with 10).
Because all subjects heard the same audiobook, we instead fit a linear ridge map from the new subject's sensors to
the training subject's sensors on stimulus-locked windows of the shared audio, using no word labels and no decoder
training. This gives +4.4 points with 5 minutes (24/28 subjects improved) and +6.5 with 10 minutes, significantly
better than supervised fine-tuning at 5 minutes (+1.7 [0.5, 2.9]). Controls show that the gain requires
stimulus-locked correspondence: the same map fitted on randomly paired windows lowers accuracy (−2.9), and matching
sensor covariances alone has no effect. Combined with fine-tuning, 10 minutes reach 32.0\%. We also report negative
results (repeated observations, more short-recording training subjects, label-free normalisation) and an
independent replication of the timing shortcut. Absolute accuracy remains low; we discuss what this implies for
calibration-light speech interfaces.
\end{abstract}

\section{Introduction}
% TODO: motivation (ALS, locked-in), gap: new subjects + no timing shortcut + minutes of data.
% Contributions: (1) replication of the timing shortcut with top-10 @50/@250 and a shift vs random-time control;
% (2) a seeded, text-disjoint calibration curve for a timing-free decoder on 28 unseen subjects;
% (3) label-free stimulus-locked alignment beats supervised calibration at 5 min, with pairing and covariance controls;
% (4) documented negative results; (5) MEG-XL fine-tuning on two 16 GB GPUs (appendix).

\section{Related work}
% Contextual word decoding \cite{dascoli2025} [to verify]; MEG-XL \cite{megxl}; timing shortcut \cite{jayalath2026shortcuts};
% LibriBrain100 and the PNPL 2026 competition \cite{libribrain100,pnpl2026};
% functional alignment / hyperalignment: fMRI decoder reuse via ridge to a reference subject \cite{ferrante2024} [abstract read],
% Haxby 2011, Chen 2015 (SRM), Thual 2023, Heo 2026 (ECoG) [to verify];
% public LibriBrain null result of the same recipe (rohit-ravi2/libribrain-2026, P1b) — must be discussed.

\section{Data and protocol}
% sub-0 training text (Sherlock 2–7, 454k words) disjoint from the test book (Sherlock 1); validation subjects
% 19–22; test subjects 1–18, 23–32; calibration on recording A, evaluation on recording B (5-gram overlap 0.6 %).
% Metric: balanced top-10 over the 50/250 most frequent training words; 3 seeds; bootstrap CIs over subjects.

\section{Method}
% Decoder: 3.6 M-param CNN on one 3 s window, contrastive to frozen T5-large word embeddings.
% Alignment: W = C_{yx} (C_{xx} + \lambda \bar\sigma^2 I)^{-1} on windows paired through the shared stimulus.

\section{Results}
\begin{table}[h]\centering
\begin{tabular}{lcc}\toprule
Gain vs zero-shot, 50 words & 5 min (28 subj.) & 10 min (18 subj.)\\\midrule
Ridge alignment (no labels) & +4.41 [+2.74, +6.03] & +6.52 [+4.77, +8.30]\\
Random pairing (control) & $-$2.88 [$-$4.10, $-$1.73] & $-$1.58 [$-$3.27, +0.03]\\
Supervised fine-tune & +2.70 [+1.80, +3.58] & +6.30 [+4.90, +7.73]\\
Alignment + fine-tune & +2.26 [+0.61, +3.85] & +7.48 [+5.51, +9.33]\\\bottomrule
\end{tabular}
\caption{28 never-seen subjects, mean of 3 seeds, 95\% bootstrap CI over subjects (results/phase3\_wide\_v1).}
\end{table}
% TODO: covariance-only and Procrustes controls (phase3_alignment_v2, phase3_weak_v1); calibration curve to 20 min;
% shortcut replication table (phase2); negative results; more-data / ensemble run (phase3_data_v1, pending).

\section{Limitations and ethics}
% Perceived speech only; same audio required for alignment; one training subject; one book; low absolute accuracy;
% LibriBrain CC-BY-NC-4.0; consent; no mind-reading claims.

\section*{Acknowledgments and AI usage}
% Honest statement: experiments, code and drafting were carried out with an AI coding assistant under the author's
% direction; all numbers are reproducible from the public repository. Adapt to the venue's policy.

\begin{thebibliography}{9}
\bibitem{jayalath2026shortcuts} D.~Jayalath, O.~Parker Jones. Removing Timing Shortcuts Improves Non-Invasive Brain-to-Text. arXiv:2609.40359, 2026.
\bibitem{megxl} D.~Jayalath, O.~Parker Jones. MEG-XL: Data-Efficient Brain-to-Text via Long-Context Pre-Training. ICML 2026, arXiv:2602.02494.
\bibitem{libribrain100} F.~Mantegna et al. LibriBrain100. arXiv:2608.25204, 2026.
\bibitem{pnpl2026} The 2026 PNPL competition. arXiv:2609.03231, 2026.
\bibitem{ferrante2024} M.~Ferrante et al. Through their eyes: multi-subject brain decoding with simple alignment techniques. arXiv:2309.00627 [to verify].
\bibitem{dascoli2025} S.~d'Ascoli et al. Decoding individual words from non-invasive brain recordings. Nat.\ Commun.\ 2025 [to verify].
\end{thebibliography}
\end{document}
